% !TeX program = xelatex
% UTF-8; compile with XeLaTeX from this directory or the repository root.
% Define \NeonatalSeizuresInNotebook before inputting into a XePersian document.
% Source: Google Drive Notebook/neonatal seizures, in filename order.
% pic01.jpg: 20261016_111444780.jpg (1Z7HOHhLLi6Vm1d0bMgToUVBQiNR5rpHw)
% pic02.jpg: 20261016_111451276.jpg (16-6trWU04M2pAMGXi_YDCWfrbwfaqvpi)
% pic03.jpg: 20261016_111454838.jpg (1eB-9gbzaESsjlpD23rq0n99izVi4FeJq)
% pic04.jpg: 20261016_111500397.jpg (1jQ5K6lQsM-S15cm6nQEhWUWJ-IypRAE3)
% pic05.jpg: 20261016_111505020.jpg (1SYe4SEEufRgovz-zqPDiG9H_o0yVUiaV)
% Arrows are rendered as subordinate statements; highlighted text is bold.
% Source spellings and doses are transcribed, not medically revised.
% Crossed-out words are explicitly recorded. Original photographs at the end
% preserve printed ruling numbers 1--9, highlighting, arrows and all marks.
\ifdefined\NeonatalSeizuresInNotebook
  \let\NSfinish\relax
\else
  \def\NSfinish{\end{document}}
  \documentclass[12pt,a4paper]{article}
  \usepackage[margin=20mm]{geometry}
  \usepackage{amsmath,amssymb}
  \usepackage{graphicx}
  \usepackage{enumitem}
  \usepackage{xepersian}
  \IfFontExistsTF{Amiri}{\settextfont{Amiri}}{\settextfont{DejaVu Sans}}
  \setlatintextfont{DejaVu Serif}
  \setlength{\parindent}{0pt}
  \setlength{\parskip}{4pt}
  \setlength{\emergencystretch}{3em}
  \setlist[itemize]{nosep,topsep=3pt}
  \begin{document}
\fi
\newcommand{\NSheader}[1]{%
  \noindent\textbf{تصویر #1 ـ صفحهٔ دست‌نویس #1}\par
  {\footnotesize\lr{Subject:}\quad ۹۴\par
  \lr{Year: \rule{8mm}{0.2pt}\quad Month: \rule{8mm}{0.2pt}\quad Date: \rule{8mm}{0.2pt}}\par}}

\section*{۹۴ ـ \lr{Hypoxic-Ischemia Encephalopathy, ICH and Seizures}}
\NSheader{۱}
\begin{itemize}
\item نوزاد بیشتر خواب است. (خواب فعال یا \lr{REM})
\item رفلکس‌ها نوزادی:
  \begin{itemize}
  \item \lr{Moro, grasp, stepping, rooting, sucking, CER}
  \item \lr{(Crossed extensor reflexes)}
  \end{itemize}
\end{itemize}
\subsection*{تشنج‌ها نوزادی}
\begin{itemize}
\item تشنج پست آسفیکسی: علت انسفالوپاتی هایپوکسیک ـ ایسکمیک.
  \begin{itemize}
  \item \textbf{۱۲–۲۴ ساعت پس از آسفیکسی}.
  \item با دوزهای عادی داروهای ضدتشنج عود می‌کند.
  \item ممکن است به علت \textbf{هایپوگلیسمی و هایپوکلسمی} ناشی از آسفیکسی باشد.
  \end{itemize}
\item تشنج به دنبال \lr{IVH}: علت شایع تشنج در نوزادان نارس.
  \begin{itemize}
  \item \textbf{در روز ۱–۳ بعد از تولد}.
  \item همراه با فونتانل بالج، \lr{CSF} خونی، آنمی، لتارژی، کما.
  \end{itemize}
\item تشنج به دنبال هایپوگلیسمی:
  \begin{itemize}
  \item \textbf{۱–۲ ساعت پس از تولد}.
  \item \textbf{۲۴–۴۸ ساعت پس از تغذیهٔ ناکافی}.
  \end{itemize}
\item تشنج به دنبال هایپوکلسمی یا هایپومنیزیمی: در نوزادان \lr{high risk}.
  \begin{itemize}
  \item به‌خوبی به درمان با کلسیم و منیزیم جواب می‌دهد.
  \end{itemize}
\item تشنج نوزاد در اتاق زایمان:
  \begin{itemize}
  \item معمولاً ناشی از تزریق مادهٔ بی‌حسی به اسکالپ نوزاد.
  \item همراه با \textbf{برادی‌کاردی نوزاد}، \textbf{مردمک دیلاته و فیکس}.
  \item به دلیل آنوکسی شدید.
  \item مالفورماسیون‌های مغزی مادرزادی.
  \end{itemize}
\end{itemize}

\clearpage
\section*{۹۴ ـ تشنج نوزادی (ادامه)}
\NSheader{۲}
\begin{itemize}
\item تشنج پس از ۵ روز اول: \textbf{عفونت، ویتدراوال دارو}.
\item تشنج همراه با لتارژی، اسیدوز، سابقهٔ خانوادگی مرگ نوزاد:
  \begin{itemize}
  \item \lr{Inborn Error of Metabolism}.
  \end{itemize}
\item نوزادی که والدینش سابقهٔ تشنج نوزادی داشته‌اند: در ریسک
  \lr{Benign familial neonatal seizures}.
  \begin{itemize}
  \item در \textbf{روز ۳–۷} شروع می‌شود و در \textbf{ماه ۱–۴} رفع می‌گردد.
  \item موتاسیون ژن \lr{KCNQ2} و \lr{KCNQ3}: کانال‌های پتاسیم.
  \end{itemize}
\item \lr{SAH}: \textbf{تشنج در نوزادی که مشکلی ندارد}، \textbf{در روز ۱–۳ تولد}، طول مدت تشنج کوتاه است، \textbf{تکرار نمی‌شود}.
\item تشنج‌های فوکال معمولاً به علت \textbf{انفارکت‌های فوکال} هستند.
\end{itemize}
\subsection*{افتراق تشنج نوزاد}
\begin{itemize}
\item \lr{Benign jitteriness}؛ \lr{Tremulousness}:
  \begin{itemize}
  \item \lr{Infant with DM mother}.
  \item \lr{Narcotic withdrawal Syndrome}.
  \item هر نوزادی پس از اپیزود آسفیکسی.
  \end{itemize}
\item جیترینس و ترمور به خلاف تشنج:
  \begin{itemize}
  \item \lr{Sensory dependent}.
  \item \textbf{\lr{Elicited by stimuli}}.
  \item \textbf{با نگه‌داشتن اندام‌ها قطع می‌شود}.
  \end{itemize}
\item جیترینس: \textbf{\lr{Rapid}} (\lr{Slow} در دست‌نویس خط خورده است)، \textbf{\lr{Fine}}.
\item تشنج: \textbf{\lr{Fast–Slow}}، \textbf{\lr{Coarse}}.
  \begin{itemize}
  \item \textbf{همراهی با حرکات غیرطبیعی چشم}.
  \item \lr{Tonic deviation to one side}.
  \end{itemize}
\item در تشنج در نوزاد با [خوانش نامطمئن: میلین کم] ممکن است حرکات \lr{tonic-clonic} دیده نشوند.
\end{itemize}

\clearpage
\section*{۹۴ ـ علائم، بررسی و درمان تشنج نوزاد}
\NSheader{۳}
\begin{itemize}
\item علائم تشنج در نوزادان: آپنه، انحراف چشم‌ها، بیرون آمدن زبان، پلک زدن، نوسان علائم حیاتی، خیره شدن.
\item اگر به تشخیص تشنج شک داشتیم: \lr{EEG}.
\item اپروچ به نوزاد با تشنج:
  \begin{itemize}
  \item \textbf{گلوکوز خون} را اندازه‌گیری کنیم؛ \lr{chemstrip}.
  \item سطح خونی سدیم، کلسیم، گلوکوز، بیلی‌روبین.
  \item اگر شک به عفونت: \textbf{کشت خون و \lr{CSF}}؛ \textbf{\lr{LP}}.
  \item خط دوم بررسی‌ها: سونوگرافی، \lr{CT}، \lr{MRI}.
  \item نقایص متابولیک را بررسی می‌کنیم.
  \end{itemize}
\item به کدام نوزادان به نقص متابولیک شک می‌کنیم؟
  \begin{itemize}
  \item لتارژی، کما، اسیدوز، کتونوری، آلکالوز تنفسی.
  \end{itemize}
\item درمان تشنج نوزاد:
  \begin{itemize}
  \item رفع علت زمینه‌ای.
  \item \textbf{فنوبارب: \lr{20–40 mg/kg}}.
  \item \textbf{فنی‌توئین: \lr{10–20 mg/kg}}.
  \item \textbf{دیازپام: \lr{0.1–0.3 mg/kg}}.
  \item درمان آلترناتیو: لوتیراستام.
  \end{itemize}
\item درمان استاتوس اپیلپتیکوس: دوزهای متعدد \textbf{فنوبارب} + \textbf{دیازپام/میدازولام}.
\end{itemize}
\subsection*{\lr{ICH}}
\begin{itemize}
\item \lr{SDH}: \textbf{ترومای حین تولد}، سفالوپلویک دیس‌پروپورشن، فورسپس، وکیوم، ماکروزومی.
  \begin{itemize}
  \item ممکن است علائم بلافاصله ظاهر نشوند.
  \item آنمی، استفراغ، تشنج، ماکروسفالی: در یک شیرخوار \textbf{۱–۲ ماهه}.
  \item \textbf{اگر علامت‌دار شد: جراحی و تخلیهٔ هماتوم}.
  \end{itemize}
\end{itemize}

\clearpage
\section*{۹۴ ـ \lr{ICH} و انسفالوپاتی هایپوکسیک ـ ایسکمیک}
\NSheader{۴}
\begin{itemize}
\item \lr{SAH}: خودبه‌خودی، در همراهی با هایپوکسی، \lr{AVM}.
  \begin{itemize}
  \item \textbf{تظاهر شایع: تشنج}.
  \item تشنج و هیدروسفالی (در صورت وجود) را درمان می‌کنیم.
  \end{itemize}
\item \lr{IVH}: در شیرخواران با \textbf{وزن هنگام تولد خیلی کم} دیده می‌شود؛ خصوصاً نارس.
  \begin{itemize}
  \item در دست‌نویس، «نارس» بالای عبارت خط‌خوردهٔ \lr{VLBW} آمده است.
  \item پاتوژنز ناشناخته است.
  \end{itemize}
\item مواردی که ریسک \lr{IVH} را بالا می‌برند:
  \begin{itemize}
  \item آسفیکسی، پنوموتوراکس، ونتیلاسیون مکانیکی، هایپرکاپنه، هایپوکسمی، لیبر طول‌کشیده، زایمان بریچ، \lr{PDA}، \lr{HF}، درمان با مایعات هایپرتونیک مثل سدیم بی‌کربنات.
  \end{itemize}
\item اکثر \lr{IVH}ها \textbf{ظرف ۳ روز اول} تولد رخ می‌دهند.
\item علائم \lr{IVH}: \textbf{تشنج}، آپنه، \textbf{برادی‌کاردی}، لتارژی، کما، \textbf{هایپوتنشن}، اسیدوز متابولیک، آنمی، فونتانل بالج، \textbf{ماتلینگ پوست}.
  \begin{itemize}
  \item با ترانسفیوژن خون اصلاح نمی‌شود.
  \end{itemize}
\item \lr{IVH} در چهار گرید دسته‌بندی می‌شود.
\end{itemize}
\subsection*{انسفالوپاتی هایپوکسیک ـ ایسکمیک}
\begin{itemize}
\item \textbf{کاهش جریان خون} رحمی ـ جفتی، هایپوکسی پری‌ناتال، \textbf{اسیدوز لاکتیک}، ایست قلبی، ایسکمی.
  \begin{itemize}
  \item \lr{local acidosis} (پیکان از اسیدوز لاکتیک).
  \end{itemize}
\item \textbf{هایپوگلیسمی} یک علت شایع است.
\item منجر به نکروز سلولی و ادم اندوتلیوم عروقی مغز.
\item علائم \lr{HIE}: ادم مغزی، [خوانش نامطمئن: بلیدینگ کورتیکال]، نکروز بازال گانگلیا؛ \lr{N.N} (یادداشت بالای «گانگلیا»).
  \begin{itemize}
  \item عقب‌ماندگی ذهنی، آتروفی مغز، \lr{CP}، فلج.
  \end{itemize}
\end{itemize}

\clearpage
\section*{۹۴ ـ \lr{HIE} (ادامه)}
\NSheader{۵}
\begin{itemize}
\item در پی پیشرفت ادم مغزی، تشنج‌ها از \textbf{۱۲–۲۴ ساعت پس از تولد} شروع می‌شوند.
\item آپگار دقیقهٔ ۱۰ اگر ۰–۳ باشد: ۲۰\% مورتالیتی و ۵۰\% خطر \lr{CP}.
\item درمان \lr{HIE} در نوزاد ترم: \textbf{هایپوترمی القایی}.
\item \lr{HIE} سه استیج مختلف دارد.
\end{itemize}

% The five unaltered original images retain the source's complete visual record.
\ifdefined\NSimages\else
  \IfFileExists{images/pic01.jpg}{\def\NSimages{images/}}{%
    \def\NSimages{pediatrics/ped-neonatology/neonatal seizures/images/}}
\fi
\clearpage
\section*{تصاویر منبع}
\begin{center}
\includegraphics[width=\linewidth,height=0.87\textheight,keepaspectratio]{\NSimages pic01.jpg}
\end{center}
\clearpage
\section*{تصویر منبع ۲}
\begin{center}
\includegraphics[width=\linewidth,height=0.87\textheight,keepaspectratio]{\NSimages pic02.jpg}
\end{center}
\clearpage
\section*{تصویر منبع ۳}
\begin{center}
\includegraphics[width=\linewidth,height=0.87\textheight,keepaspectratio]{\NSimages pic03.jpg}
\end{center}
\clearpage
\section*{تصویر منبع ۴}
\begin{center}
\includegraphics[width=\linewidth,height=0.87\textheight,keepaspectratio]{\NSimages pic04.jpg}
\end{center}
\clearpage
\section*{تصویر منبع ۵}
\begin{center}
\includegraphics[width=\linewidth,height=0.87\textheight,keepaspectratio]{\NSimages pic05.jpg}
\end{center}
\NSfinish
